% Generated draft: requires mathematical and source review.
\documentclass[11pt]{article}
\usepackage[margin=25mm]{geometry}
\usepackage{fontspec}
\setmainfont{Latin Modern Roman}
\usepackage{amsmath,amssymb,mathtools,unicode-math}
\setmathfont{Latin Modern Math}
\usepackage{enumitem,needspace,xurl,hyperref}
\setlength{\parindent}{0pt}
\setlength{\parskip}{5pt}
\setlength{\emergencystretch}{3em}
\newcommand{\N}{\mathbb{N}}
\newcommand{\Z}{\mathbb{Z}}
\newcommand{\Q}{\mathbb{Q}}
\newcommand{\R}{\mathbb{R}}
\newcommand{\C}{\mathbb{C}}
\newcommand{\F}{\mathbb{F}}
\newcommand{\A}{\mathbb{A}}
\newcommand{\PP}{\mathbb{P}}
\newcommand{\E}{\mathbb{E}}
\newcommand{\eps}{\varepsilon}
\newcommand{\dd}{\,\mathrm{d}}
\begin{document}
\section*{UnsolvedMath ID: 30000822}
\textbf{Source problem code:} catalogue entry
\subsection{Definitions and mathematical statement}
The report's Theorem 3 concerns an $r$-connected space $X$ with $r\ge1$ and constructs, from a cohomology class $u\in H^*(X)$ of sufficiently large category weight, a class
\[
 \phi(u)\in H^*(X\times X)
\]
with a prescribed cohomology-weight lower bound. The restored statement further specifies that $\phi(u)$ is a zero-divisor. The question is whether an analogous construction, retaining the appropriate lower bound, exists when $X$ is not simply connected.

For a coefficient system in which the diagonal map is defined, the usual zero-divisor condition is
\[
 \delta^*\phi(u)=0,\qquad \delta:X\longrightarrow X\times X,
 \quad\delta(x)=(x,x).
\]
No K\"unneth identification of this kernel with an algebraic tensor-product expression is assumed. Category weight ordinarily records vanishing behavior under maps of small Lusternik--Schnirelmann category; the target cohomology-weight notion must be defined by the particular theorem. The source supplies neither the precise weight definitions and normalization nor the threshold for $u$ or the lower-bound formula.

The extension must therefore recover the theorem's actual construction and quantitative weight requirement rather than merely produce an arbitrary element in the diagonal kernel. Its degree, nonvanishing requirements, naturality, and dependence on coefficients are also not stated. A trivial zero class cannot be accepted as an intended construction without knowing those requirements.

The original $r\ge1$ hypothesis implies simple connectedness in its usual path-connected setting. For the proposed nonsimply connected analogue, the allowed space category, connectedness, fundamental-group hypotheses, and possible local coefficients are not specified. No replacement $r=0$, restriction on $\pi_1$, or explicit formula such as a difference of pullbacks is silently inserted. The supplied recent analogue is not assumed to establish this exact missing theorem.
\subsection{Short English statement}
The question asks for a nonsimply connected extension of a specific high-category-weight construction, including its zero-divisor output and weight lower bound. The theorem's formula, thresholds, coefficients, and nontriviality conditions are missing, so producing any zero-divisor would not suffice.
\subsection{Sources}
\begin{itemize}
\item Local UnsolvedMath registry, permanent ID 30000822, code catalogue entry, ``Category-Weight Constructions for Nonsimply Connected Spaces.'' Proposed year 2007; proposer unspecified.
\item \emph{Mini-Workshop: Topology of closed one-forms and Cohomology Jumping Loci} (2008), Oberwolfach Reports, pp.\ 2321--2360, DOI 10.4171/owr/2007/40; \url{https://doi.org/10.4171/owr/2007/40}. The displayed report year 2008 differs from the DOI and proposed year 2007.
\item Supplied ``recent closed-form analogue'': \url{https://arxiv.org/abs/2606.09184}.
\item The retained original identifies the simply connected result as Theorem 3; its hypotheses, construction, and numerical bound are not included.
\item Statement status corrected\_verified; the supplied note says official-report content was checked and wording/notation normalized without changing the source claim.
\item Literature triage checked 2026-08-21: open. It says the report gives the simply connected result but not the genuine nonsimply connected extension. Exact theorem data and current status require independent verification. Created 2026-08-25; updated 2026-09-16.
\item Registry curation by Ulam AI / UnsolvedMath Contributors, CC BY 4.0; underlying report and literature provenance remain as supplied. Unreviewed statement-only draft.
\end{itemize}
\end{document}
